\documentclass[10pt,a4paper]{amsart}
\usepackage[margin=19mm]{geometry}
\usepackage{amsmath,amssymb,mathtools}
\usepackage[scr=rsfs,cal=euler]{mathalfa}
\usepackage{xfrac}
\usepackage[unicode,hidelinks,bookmarks=false]{hyperref}
\newcommand{\ie}{i.e.\ }
\newcommand{\cf}{cf.\ }
\newcommand{\resp}{respectively\ }
\newcommand{\Subsection}{\subsection}
\providecommand{\nobreakdash}{\nobreak\hbox{-}}
\setlength{\emergencystretch}{2em}
\multlinegap0pt
\usepackage{xcolor}
\usepackage{amssymb}
\usepackage{mathtools}
\newtheorem{definition}{Definition}
\newtheorem{lemma}{Lemma}
   \renewcommand{\P}{\mbox{${\mathcal P}$}}
\title{A Non-commutative Individual Ergodic Theorem Along Sparse Random Subsequences}
\author{L\'eonard Cadilhac, Christian Le Merdy, Safoura Zadeh}
\date{Source: arXiv:2609.07970v1 (CC BY 4.0)}
\begin{document}
\maketitle

\noindent\emph{Source and licence.} A Non-commutative Individual Ergodic Theorem Along Sparse Random Subsequences. Version arXiv:2609.07970v1 (CC BY 4.0), distributed by arXiv under the Creative Commons Attribution 4.0 International licence, http://creativecommons.org/licenses/by/4.0/, as recorded in the arXiv licence metadata for this exact version and shown on the abstract page https://arxiv.org/abs/2609.07970v1. Full licence notice: \texttt{licenses/arxiv-2609.07970-CC-BY-4.0.txt}.\par\smallskip

\noindent\textit{Editorial scope.} Counted unit: the article's lemma 3Y1 (source0.tex lines 580--599). The article's complete original proof of it is delivered with it (source0.tex lines 601--681, one \texttt{proof} environment of 81 lines). No mathematics is written, repaired or re-derived: the statement and the proof are the article's own lines, in the article's own notation, and no retained line is substituted or re-flowed.  Retained uncounted as the source-local context the proof needs: lines 456--459; lines 560--572.  Nothing was omitted from any retained span, so the package carries no omission record.  No retained line contains a citation, so the wrapper carries no bibliography and no bibliography entry.  No retained line contains a figure environment, an image-inclusion command or drawing code, so no image is delivered and the package declares no asset dependency.  Editorial formatting only, in addition: an AMS \texttt{amsart} preamble that restates the article's own declarations and macros used by the retained lines, namely the environments \texttt{definition}, \texttt{lemma} on separate counters, together with the standard packages those lines need.  The retained environments print in this document as Definition 1, Lemma 1 in order of appearance, whereas the article prints the same items with the numbering of its own full document.  The complete native source of the article as distributed in the arXiv e-print for this exact version is bundled unchanged as \texttt{sources/209680/source0.tex}.
\begin{equation}\label{2Tr}
\tau\bigl((e\wedge f)^\perp\bigr)
\leq \tau(e^\perp)+\tau(f^\perp),
\end{equation}

\begin{definition}\label{3MI}
Let \((V_j)_{j\ge1}\) be a sequence of bounded 
maps \(V_j : L^1(M) \longrightarrow L^1(M)\).
We say that \((V_j)_{j\ge1}\) satisfies a weak-type
$(1,1)$-maximal inequality if there is a constant \(C\geq 0\) 
such that for every \(x\in L^1(M)\) and every \(\lambda>0\), 
there exists a projection \(e\in \P(M)\) satisfying
\begin{equation}\label{3Weak11}
\tau(e^\perp)\le \frac{C\|x\|_1}{\lambda}
\qquad\text{and}\qquad
\|eV_j(x)e\|_\infty \le \lambda, \quad \forall j\ge1.
\end{equation}
\end{definition}

\begin{lemma}\label{3Y1}
Let $(V_j)_{j\geq1}$ be a sequence of bounded maps on $L^1(M)$ that
satisfies a weak-type $(1,1)$-maximal inequality. Let $x\in L^1(M)$,
and suppose that there exists a sequence $(y_m)_{m\geq1}$ in $L^1(M)$
such that
\[
\|x-y_m\|_1\longrightarrow0
\qquad\text{as }m\to\infty,
\]
and for every $m\geq1$,
\[
V_j(y_m)\longrightarrow0
\qquad\text{b.a.u.\ as }j\to\infty.
\]
Then
\[
V_j(x)\longrightarrow0
\qquad\text{b.a.u.\ as }j\to\infty.
\]
\end{lemma}

\begin{proof}
Let $C\geq0$ be a constant for the weak-type $(1,1)$-maximal inequality in Definition \ref{3MI}, and fix $\varepsilon>0$. Passing to a subsequence, 
we may assume that
$\|x-y_m\|_1\leq4^{-m}\varepsilon$,
for all $m\geq1$.
Set $z_m=x-y_m$, so that $x=y_m+z_m$ and
\[
\|z_m\|_1\leq4^{-m}\varepsilon,
\qquad m\geq1.
\]
Applying the weak-type $(1,1)$-maximal inequality to $z_m$ with $\lambda=2^{-m}$, we obtain projections $e_m\in\P(M)$ such that
\(\tau(e_m^\perp)\leq C\varepsilon 2^{-m}\)
and
\[
\|e_mV_j(z_m)e_m\|_\infty\leq2^{-m},
\qquad j\geq1.
\]
Let
\(e=\bigwedge_{m=1}^\infty e_m.
\)
By \eqref{2Tr},
\[
\tau(e^\perp)
\leq\sum_{m=1}^\infty\tau(e_m^\perp)
\leq C\varepsilon.
\]

On the other hand, since for every $m\geq1$, $V_j(y_m)\to0$ b.a.u.\ as $j\to\infty$, for each $m\geq1$ there exists a projection $f_m\in\P(M)$ such that
\[
\tau(f_m^\perp)\leq\varepsilon2^{-m}
\]
and
\[
\|f_mV_j(y_m)f_m\|_\infty\longrightarrow0
\qquad\text{as }j\to\infty.
\]
Set
\(f=\bigwedge_{m=1}^\infty f_m
\ \text{and}\
g=e\wedge f.\)
Again by \eqref{2Tr},
\(\tau(f^\perp)
\leq\sum_{m=1}^\infty\tau(f_m^\perp)
\leq\varepsilon, \)
and hence
\[
\tau(g^\perp)
\leq\tau(e^\perp)+\tau(f^\perp)
\leq(1+C)\varepsilon.
\]

It remains to show that
\(\|gV_j(x)g\|_\infty\longrightarrow0
\ \text{as }j\to\infty.\)
Fix $m\geq1$. Since $x=y_m+z_m$, the triangle inequality gives
\[
\|gV_j(x)g\|_\infty
\leq
\|gV_j(y_m)g\|_\infty
+
\|gV_j(z_m)g\|_\infty.
\]
Since $g\leq f_m$ and $g\leq e_m$, we obtain
\[
\begin{aligned}
\|gV_j(x)g\|_\infty
&\leq
\|f_mV_j(y_m)f_m\|_\infty
+
\|e_mV_j(z_m)e_m\|_\infty\\
&\leq
\|f_mV_j(y_m)f_m\|_\infty+2^{-m}.
\end{aligned}
\]
Therefore,
\[
\limsup_{j\to\infty}\|gV_j(x)g\|_\infty
\leq2^{-m}.
\]
Since $m\geq1$ was arbitrary, letting $m\to\infty$ yields \(\lim_{j\to\infty}\|gV_j(x)g\|_\infty = 0\), which completes the proof.
\end{proof}

\end{document}
